% ============================================================
% Features and Labels Table — Hieroglyph Analogy
% Supervised Learning — CSC491 Textbook, Chapter: Machine Learning
%
% Follows directly after tab:hieroglyph-labeled.
% Shows how raw inputs are decomposed into measurable features,
% and paired with a label the model is trained to predict.
%
% Preamble requirements:
%     \usepackage{fontspec}
%     \newfontfamily\egyptfont{Noto Sans Egyptian Hieroglyphs}
%     \usepackage{booktabs, array, xcolor}
% ============================================================

\begin{table}[ht]
\centering
\setlength{\tabcolsep}{10pt}
\renewcommand{\arraystretch}{1.8}
\caption{%
  Features and labels extracted from a labeled hieroglyph dataset.
  Rather than feeding the raw image directly to the model, we first
  describe each glyph in terms of measurable \textbf{features} --- the
  properties the model can reason about. The \textbf{label} is the
  correct answer we want the model to learn to predict.%
}
\label{tab:hieroglyph-features}
\begin{tabular}{
  >{\centering\arraybackslash}p{1.6cm}   % glyph
  >{\centering\arraybackslash}p{1.6cm}   % curved strokes
  >{\centering\arraybackslash}p{1.6cm}   % straight strokes
  >{\centering\arraybackslash}p{1.8cm}   % has eyes?
  >{\centering\arraybackslash}p{1.8cm}   % animal?
  >{\centering\arraybackslash}p{1.8cm}   % vertical?
  !{\color{black!25}\vrule width 0.8pt}  % divider
  >{\centering\arraybackslash}p{2.0cm}   % label
}
\toprule
\multicolumn{6}{c}{\textbf{Features} $\mathbf{x}$}
  & \multicolumn{1}{c}{\textbf{Label} $\mathbf{y}$} \\
\cmidrule(r){1-6} \cmidrule(l){7-7}
\textbf{Glyph}
  & \textbf{Curved strokes}
  & \textbf{Straight strokes}
  & \textbf{Has eye}
  & \textbf{Animal}
  & \textbf{Vertical}
  & \textbf{Translation} \\
\midrule
{\egyptfont 𓂀} & 5 & 2 & Yes & No  & No  & Eye of Horus \\
{\egyptfont 𓅓} & 3 & 4 & Yes & Yes & No  & Owl ($m$)    \\
{\egyptfont 𓈖} & 6 & 0 & No  & No  & No  & Water ($n$)  \\
{\egyptfont 𓆑} & 4 & 3 & No  & Yes & No  & Viper ($f$)  \\
{\egyptfont 𓏏} & 1 & 3 & No  & No  & No  & Bread ($t$)  \\
{\egyptfont 𓃭} & 2 & 5 & Yes & Yes & No  & Lion ($r$)   \\
{\egyptfont 𓇋} & 0 & 2 & No  & No  & Yes & Reed ($i$)   \\
{\egyptfont 𓎛} & 2 & 6 & No  & No  & No  & Rope ($h$)   \\
\bottomrule
\end{tabular}

\smallskip
\footnotesize
\textit{Note:} These features are simplified for illustration.
In practice, a computer vision model would extract hundreds of such
measurements automatically, but the principle is identical.
Each row is one training example; the model's job is to learn which
combinations of feature values reliably predict each label.

\end{table}